\paragraph{My contribution.} I helped set up the project plan, split the work into three streams with shared contracts between them, and scaffolded the repository. I then led Stream C, the reasoning and cognitive layer. I built the rule engine, which uses forward chaining, along with its validator, guardrails and rule versioning in the app. I built the Neo4j knowledge graph and the cohort signals it feeds back into the rules. I also built the staff assistant, which routes a question to a tool and answers only from that tool's output. Beyond that, I built the app database with migrations, seeds and login by role, ran experiments E6 to E8, and did general UI/UX work in the web app.

\paragraph{What I learned.} The most useful lesson was how the three course modules fit together in one system. From Machine Reasoning, writing rules as data with a validator taught me that a rule base needs the same care as code. It needs tests that state which rules must fire, checks for contradictions, and versions that can be rolled back. From Cognitive Systems, the knowledge graph showed me that structure can add reach rather than accuracy. It flagged 243 members at low risk whose routine had just been disrupted, members the churn model would never rank. The assistant taught me that grounding matters more than the model. Its answers were 95.1\% supported by evidence with both a hosted and a local model, but its intent routing was no better than keywords on unseen questions, so the weak point was the routing and not the language model.

The hardest part was the trade off between guardrails and plan value. Guardrails removed 31.3\% of proposals, almost all of them repeat contacts within 30 days. When staff time was tight this cost almost nothing, but at 120 staff hours the plans without guardrails kept 6 to 14 more members, because our scenario has no contact fatigue effect to penalise repeat contacts. I learned to be honest about what an evaluation can measure, so we claim safety and auditability for the guardrails and not a retention gain. Next time I would build the cost of the behaviour a rule protects against into the evaluation from the start.

Integration was the other big lesson. Machine learning, rules, case based reasoning and optimisation each worked alone well before they worked together. Agreeing on schemas early, and keeping the system running when the LLM or the graph was unavailable, saved us more time than any single model did.

\paragraph{Applying it.} I work as a product development engineer at LAB7, the venture building arm of Aramco, where many of the startups we build and support are developing AI systems. This project gives me a practical checklist for them. Keep business rules explicit and versioned beside the model, ground any LLM feature in tool output and measure its faithfulness, use a knowledge graph where relationships carry signal, and decide early how a system will be evaluated when real outcomes are not yet available. Above all, I now push for explainability as a design requirement, since staff will only act on a plan they can question.
